\documentclass[10pt,a4paper]{amsart}
\usepackage{amsmath,amssymb,amsthm,graphicx}
\usepackage{amsmath}
\usepackage{amssymb}
\usepackage{hyperref}
% Minimal class-compatibility preamble for a reader excerpt of
%   Maxim Nazarov, "Yangian of the General Linear Lie Superalgebra",
%   SIGMA 16 (2020), 112, DOI 10.3842/SIGMA.2020.112.
%
% The block below is the author's own macro and theorem declaration block of the pinned
% source file, copied mechanically and unchanged (source lines 9-71). The SIGMA publisher
% class sigma.cls is not installed here, so the wrapper class declares the same theorem
% environments; no mathematical macro, statement, symbol or assumption is introduced,
% changed or removed by this file.
\numberwithin{equation}{section}

\newtheorem{Theorem}{Theorem}[section]
\newtheorem{Corollary}[Theorem]{Corollary}
\newtheorem{Lemma}[Theorem]{Lemma}
\newtheorem{Proposition}[Theorem]{Proposition}

\def\1{^{-1}}

\def\a{\mathfrak{a}}
\def\al{\alpha}

\def\be{\beta}
\def\bi{{\bar\imath }}
\def\bj{{\bar\jmath }}
\def\bk{{\bar k }}
\def\bl{{\bar l }}
\def\bp{{\bar p }}
\def\bq{{\bar q }}
\def\br{{\bar r }}
\def\bs{{\bar s }}

\def\CC{{\mathbb{C}}}
\def\CMN{\CC^{M|N}}

\def\de{\delta}
\def\De{\Delta}

\def\End{\operatorname{End}}
\def\ep{\varepsilon}

\def\g{\glMN [u]}
\def\ge{\geqslant}
\def\glMN{\mathfrak{gl}_{M|N}}
\def\gr{{\rm gr} }
\def\grpr{{\rm gr}^{\prime} }

\def\id{{\mathrm{id}}}

\def\le{\leqslant}

\def\ns{{\hskip-.5pt}}

\def\om{\omega}
\def\ot{\otimes}

\def\RS{R^{\sharp}}

\def\S{\operatorname{S}}
\def\si{\sigma}
\def\str{\operatorname{str}}
\def\Sym{\mathfrak{S}}

\def\TD{\dot{T}}
\def\TP{T^{\prime}}
\def\TS{T^{\sharp}}

\def\Ua{\operatorname{U}(\a [u])}
\def\UMN{\operatorname{U}(\glMN)}
\def\Ug{\operatorname{U}(\g)}
\def\XP{X^{\prime}}
\def\YMN{\operatorname{Y}(\glMN)}
\def\ZZ{\mathbb{Z}}
\title{Triviality of the centre of a Lie superalgebra passes to the enveloping algebra of its polynomial current algebra}
\author{Maxim Nazarov}
\date{doi:10.3842/SIGMA.2020.112; CC BY 4.0}
\begin{document}
\maketitle
\noindent\textit{Editorial note.} All delivered mathematical text is the author's own, copied from the pinned publisher source sources/algebra-fl05-yangian-of-the-general-linear-lie/pinned-original-source.tex; no statement or proof step was rewritten, reordered or removed.\par\medskip
\noindent\textit{Editorial note.} Preamble provenance: the retained preamble is the source's own retained-preamble.tex, whose header records that its body is the author's macro and theorem declaration block of pinned-original-source.tex lines 9-71 copied mechanically and unchanged; the wrapper class supplies the amsmath baseline that this author block assumes (the SIGMA publisher class is not installed, and the pinned file itself opens with a RequirePackage/ifpdf/documentclass block that must not be re-injected as macro text). Nothing was added, no macro was invented, and the excerpt serves only as preamble, never as a span coordinate.\par\medskip
\noindent\textit{Editorial note.} Assembly history: r001 and r002 both failed with 'Undefined control sequence \textbackslash{}numberwithin' because the card carried the author macro block twice - once through retained-preamble.tex and once as a second span over pinned lines 2-71, which also pulled in the source's own RequirePackage/ifpdf/documentclass block and suppressed the wrapper class's baseline packages. r003 keeps a single preamble span (retained-preamble.tex 1-72). No mathematical text, span or statement changed between the revisions.\par\medskip
\noindent\textit{Editorial note.} External references are navigation pointers only; every deliberately kept step of the author's argument stays inside the counted spans. The displayed numbers of those pointers are the author's own numbering under the author's \textbackslash{}newtheorem declarations (the source has no subsections, so numbering is section.index), computed mechanically by the lane's scanner.\par\medskip
\noindent\textit{Editorial note.} The proof refers only to equations labelled inside the counted span (A, hq, ad) and to one bibliography item, so no navigation stub was needed and the built document carries no undefined reference.\par\medskip
\noindent\textit{Editorial note.} This is the first card of this lane from this source; the source was picked after the licence check (license-observation.json present, status supported) and after the same-paper twin check came out negative.\par\medskip
\noindent\textit{Allowed standard prerequisites (editorial).}\begin{itemize}
\item Reduction of the centre of the enveloping algebra to the invariants of the adjoint action, via the Poincare-Birkhoff-Witt theorem. Cited inside the counted proof at lines 1247-1252. Shows that triviality of the centre of U(a[u]) follows once the invariants of the adjoint action of a[u] on its supersymmetric algebra are trivial.
\item A homogeneous basis of the Lie superalgebra and its structure constants. Used inside the counted proof at lines 1255-1261. Provides the coordinates f\_1, ..., f\_K, their Z\_2-degrees and the constants C\_\{pq\}\textasciicircum{}r in which the whole computation is written.
\end{itemize}\medskip
\begin{Proposition}
\label{L27}
Suppose the centre of the Lie superalgebra $\a$ is trivial.
Then the centre of the universal enveloping algebra $\Ua$
is also trivial, that is equal to $\CC $.
\end{Proposition}

\begin{proof}
Consider adjoint action of the Lie superalgebra $\a [u]$
on its supersymmetric algebra. By
Poincar\'e--Birkhoff--Witt theorem for $\Ua$
it suffices to show that the space
of invariants of~this action is trivial,
see \cite[Theorem~5.15]{MM}.

Let $A$ be an element of the supersymmetric algebra of $\a [u]$
invariant under the adjoint action. Let $K=\dim\a $.
Choose any $\ZZ_2 $-homogeneous basis $f_1,\dots,f_K$ of $\a$ and
write the Lie brackets~as
\begin{gather*}
[ f_p ,f_q ]=\sum_{r=1}^K C_{pq}^{r} f_r,
\end{gather*}
where $C_{pq}^{r}\in\CC$ are the structure constants.
Let $d$ be the maximal degree such that $A$ depends on at least one
of the basis
elements $f_1 u^{d},\dots, f_K u^{d}$ of $\a [u] $. Write
\begin{gather}
\label{A}
A =\sum_{m_1,\dots,m_K}A_{m_1\dots m_K}
\big( f_1 u^{d} \big)^{m_1}\cdots\big( f_K u^{d} \big)^{m_K},
\end{gather}
where the coefficients $A_{m_1\dots m_K}$ are certain elements
of the supersymmetric algebra which do~not depend
on $f_1 u^{d},\dots, f_K u^{d} $.
For each $p=1,\dots, K$ let $\bp$ be the $\ZZ_2 $-degree of the basis element
$f_{p}$ of $\a $. We allow $m_1,\dots, m_K$ in \eqref{A}
to range over $0,1,2, \dots$ but assume that $A_{m_1\dots m_K}=0$
if $m_{p}>1$ for at least one index $p$ with $\bp=1 $.

For each $p=1,\dots, K$ we have the equation
$\text{ad} ( f_p u )(A)=0$ in the supersymmetric algebra
of $\a [u] $. In particular,
the component of the left-hand side of this equation
that involves any of the basis elements
$f_1 u^{d+1},\dots, f_K u^{d+1}$ must be zero. Thus
\begin{gather*}
\sum_{m_1,\dots,m_K}A_{m_1\dots m_K}
\sum_{q,r=1}^K
m_{q} C_{pq}^{r}
\big( f_1 u^{d} \big)^{m_1}
\cdots
\big( f_q u^{d} \big)^{m_{q}-1}
\cdots
\big( f_K u^{d} \big)^{m_K}
f_r u^{d+1}
(-1)^{h_{q}}=0,
\end{gather*}
where we used the notation
\begin{gather}\label{hq}
h_{q}=
\sum_{s=1}^{q-1} \bp \bs m_{s}
 + \sum_{s=q+1}^K\br \bs m_{s} .
\end{gather}
If follows that for any non-negative integers
$m_1^{\prime},\dots,m_K^{\prime}$ we have the equations
\begin{gather}\label{ad}
\sum_{q=1}^M
A_{m_1^{\prime}\dots m^{\prime}_q+1\dots m^{\prime}_K}
(m^{\prime}_q+1) C_{pq}^{r} (-1)^{h_{q}^{\prime}}=0,
\qquad\text{where}\quad
p ,r=1,\dots, K .
\end{gather}
Similarly to \eqref{hq} here we used the notation
\begin{gather*}
h_{q}^{\prime}=\sum_{s=1}^{q-1} \bp \bs m_{s}^{\prime}
 + \sum_{s=q+1}^K\br \bs m_{s}^{\prime} .
\end{gather*}

Let us now fix any non-negative integers
$m_1^{\prime},\dots,m_K^{\prime}$
and observe that the elements
\begin{gather*}
f_q^{\prime}=(m_q^{\prime}+1) f_q
 (-1)^{h_{q}^{\prime}}
\qquad\text{with}\quad q=1,\dots, K
\end{gather*}
also form a basis of $\a $. Since the centre of
the Lie superalgebra $\a$ is trivial, the system
\begin{gather*}
\left[ f_p , \sum_{q=1}^K w_{q} f_q^{\prime} \right]=0
\qquad\text{with}\quad
p=1,\dots, K
\end{gather*}
of linear equations on $w_{1},\dots, w_{K}\in\CC$
has only zero solution. This system can be written~as
\begin{gather*}
\sum_{q=1}^K w_{q}
(m^{\prime}_q+1) C_{pq}^{r} (-1)^{h_{q}^{\prime}}=0,
\qquad\text{where}\quad
p ,r=1,\dots, K .
\end{gather*}
Hence by comparing the latter system with \eqref{ad} we obtain that
$A_{m_1^{\prime}\dots m^{\prime}_q+1\dots m^{\prime}_K}=0$
for each index $q=1,\dots, K $. It follows that $A\in\CC$
and that $d=0$ in particular.
\end{proof}

% Bibliography excerpt for this card, copied verbatim from the pinned publisher source
% sources/algebra-fl05-yangian-of-the-general-linear-lie/pinned-original-source.tex,
% whose own \begin{thebibliography} block starts at line 2227. Only the entries cited
% inside this card's counted spans are reproduced, character for character, with no
% reformatting; the wrapper class supplies hyperref, which these entries use.
\begin{thebibliography}{99}

\bibitem{MM}
Milnor J.W., Moore J.C., On the structure of {H}opf algebras, \href{https://doi.org/10.2307/1970615}{\textit{Ann. of
 Math.}} \textbf{81} (1965), 211--264.

\end{thebibliography}
\end{document}
